% AJPM
\documentclass[14pt]{extbook} % [Tamaño del papel] y {Tipo del documento}

% +-----------------+
% |  Cosas básicas  |
% +-----------------+

% Encoding
\usepackage[utf8]{inputenc}			% Usar el encoding en UTF8 (ayuda para caracteres especiales)
\usepackage[T1]{fontenc}			% Más ayuda para tener caracteres especiales
\usepackage{tipa}					% Permite usar los caracteres fonéticos chéveres que sólo aparecen en la página en la que se muestra cómo pronunciar ex hypothesi

% Lenguaje
\usepackage[spanish]{babel}			% Configurar Babel para que maneje 
										% el documento según las reglas del español

% Vínculos internos
\usepackage{hyperref}			% Hace que los puntos del índice redirijan a las partes del PDF. ¿Genial, no?

% Empezar el conteo de notas (pies de página) en cada capítulo
\usepackage{etoolbox}			% No sé si esto deba ir aquí ¯\_(ツ)_/¯
\makeatletter
	\pretocmd{\@schapter}{\setcounter{footnote}{0}}{}{}
	\pretocmd{\@chapter}{\setcounter{footnote}{0}}{}{}
\makeatother

% +-----------------+


% +-----------------------------------+
% | Orden del documento y otras cosas |
% +-----------------------------------+
\renewcommand{\contentsname}{Índice}				% Cambiar título del índice
\usepackage{tocvsec2}								% Libreria para controlar mejor qué aparece en el índice
\maxtocdepth{chapter}						% Sólo muestra los capítulos en el índice, ignorando las secciones y todo lo demás
% +-----------------+


% +-----------------------------+
% | Visualización del documento |
% +-----------------------------+

% Cambiar margenes
\usepackage[letterpaper, margin=3cm, bmargin=3.5cm]{geometry} % Papel carta, todos los margenes de 3cm, mover la línea con número de página y demás info

% Cabeceras
\usepackage{fancyhdr}					% Permite modificar la cabecera y la parte inferior (ni idea cómo se llame) de la página
\pagestyle{fancy}						% Usar fancyhdr en el estilo de las páginas
\fancyhf{}								% Quitar lo que sea que tenga la cabecera y el pié de página
\renewcommand{\headrulewidth}{0pt}		% Quitar la línea superior (reducirla a 0 puntos)

% Interlineado
\usepackage{setspace}

% Usar Cormorant Garamond como tipografía
\usepackage{fontspec}
\usepackage[T1]{fontenc}
\setmainfont[Path = fonts/, BoldFont={Cormorant-bold.ttf}, ItalicFont={Cormorant-italic.ttf},BoldItalicFont={CormorantGaramond-BoldItalic.ttf}]{Cormorant-Regular.ttf}

% Pa poder usar tablas
\usepackage{longtable}
\usepackage{booktabs}
\usepackage{array}

% Usar imágenes
\usepackage{graphicx}
\graphicspath{ {img/} }				% La carpeta en la que se guardan las imágenes

% Colores
\usepackage[table]{xcolor}						% Librería para tener un mejor manejo de colores
\definecolor{ehAzul1}{rgb}{0.35,0.55,0.85}		% Azul oscuro
\definecolor{ehAzul2}{rgb}{0.75,0.85,0.95}		% Azul claro
\hypersetup{									% Configurar colores de los vínculos dentro del documento
	colorlinks=true,
	linkcolor=black,
	urlcolor=ehAzul1,
	citecolor=ehAzul1,
}

\urlstyle{same}

% +-----------------+


% +---------------+
% |  Maquetación  |
% +---------------+

% Manejo de párrafos
\setlength{\parskip}{1em}			% Espacio entre párrafos

% +-----------------+


% +-----------+
% |  Títulos  |
% +-----------+

% Librería para modificar títulos
\usepackage{titlesec}
\usepackage{titling}

% Cambiar los títulos de los capítulos
\titleformat
{\chapter} 								% Especifica que se van a editar los títulos de los capítulos
[display] 								% La configuración de cómo mostrar el título, en este caso centrado
{\linespread{0}\large\bfseries} 		% Texto centrado\negrilla\texto grande\itálica
{\thechapter}				 			% Cómo poner el número del capítulo, si se quiere
{0.5ex} 								% Separador
{\vspace{1cm}} 										% Código antes del título
[\vspace{1cm}] 										% Código después del título

\titlespacing{\chapter}{0.8cm}{-1ex}{1pc}	% Cuadrar bién el título de los artículos


% Cambiar los títulos de las secciones
\renewcommand{\thesection}{\arabic{section}}	% Modifica la numeración de las secciones para que sólo ponga el número de la sección y no el del capítulo (evita secciones 0.1, 0.2, 0.3, etc)
\titleformat
{\section}									% Especifica que se van a editar los títulos de las secciones
[hang]										% La configuración de cómo mostrar el título, en este caso justificado a la izquierda
{\normalsize\bfseries}							% Formato en el que aparecerá el texto, letra grande y en negrita (bold)
{\thesection.}								% El número de la sección
{1ex}										% Espacio entre el número y el nombre de la sección
{}											% Código antes del título
[]											% Código después del título

\titlespacing{\section}{0.8cm}{0.5cm}{0cm}	% Cuadrar bién las secciones, que por el interlineado entre párrafos queda inmundo

% Manejo de títulos de cada capítulo con su respectivo autor
% Amamos (sí, tú también) a este extraño :D https://tex.stackexchange.com/questions/156862/displaying-author-for-each-chapter-in-book
\usepackage{suffix}

% Se define el comando usado para poner el autor y el título en el índice
\newcommand\chapterauthor[1]{\authortoc{#1}\printchapterauthor{#1}}
\WithSuffix\newcommand\chapterauthor*[1]{\printchapterauthor{#1}}

% Comando para poner el autor del capítulo debajo del título
\makeatletter
\newcommand{\printchapterauthor}[2]{%
	{\parindent0.8cm\vspace*{-45pt}%
		\linespread{1.1}\small\scshape#1\footnote{#2}%
		\par\nobreak\vspace*{50pt}}
	\@afterheading%
}

% Comando para poner el autor en el índice
\newcommand{\authortoc}[1]{%
	\addtocontents{toc}{\vskip-20pt}		% Especifica el espacio entre el título del artículo y su autor
	\addtocontents{toc}{%
		\protect\contentsline{chapter}%
		{\hskip1.5em\mdseries\scshape\protect\footnotesize#1}{}{}}
	\addtocontents{toc}{\vskip-10pt}		% Especifica el espacio entre el nombre del autor y el siguiente título
}
\makeatother

% Comando para poner título, autor, información del autor, y poner título y autor en la cabecera y pié de página
\makeatother
\newcommand{\infoCap}[4]{			% Argumentos: 1) Título del artículo 1 2) Título del artículo 2 
									% 3) Autor 4) Información del autor
									% A veces es necesario poner el título en más de una línea
									% (En especial porque los filósofos ponen títulos DEMASIADO largos),
									% pero poner \\ dentro del arg #1 jode con la función que pone
									% el título en la cabecera de la página. Por eso el título está dividido
									% en dos argumentos diferentes. El segundo se puede dejar vadío si no
									% hace falta dividir el título, y no pasa nada malo. Es un workaround
									% horrible, pero funciona
	\chapter*{#1\\#2}				% Títuloombre del capítulo
	\addcontentsline{toc}{chapter}{\normalsize #1~#2}				% Añadir el título del capítulo al índice
	\chapterauthor{#3}{#4}			% Poner al autor del capítulo debajo del título y en el índice
	\fancyfoot[RE,RO]{{\footnotesize #1\\#2\\\scshape\vspace{-0.2cm}#3}}	% Poner título y autor al lado derecho del pie de página
	% Esta vaina son las cabeceras. Ni p idea de cómo funcionan :D
	\fancyhead[LE]{\hspace*{-0.2\headwidth}\colorbox{ehAzul2}{\makebox[\dimexpr0.2\headwidth-2\fboxsep][c]{\strut\textcolor{ehAzul1}{\bfseries\thepage}}}%
		\colorbox{ehAzul1}{\makebox[\dimexpr\headwidth-2\fboxsep][l]{\strut\scshape\textcolor{ehAzul2}{\bfseries#1\ldots}}}}
	\fancyhead[LO]{\colorbox{ehAzul1}{\makebox[\dimexpr\headwidth-2\fboxsep][r]{\strut\scshape\textcolor{ehAzul2}{\bfseries#1\ldots}}}%
		\colorbox{ehAzul2}{\makebox[\dimexpr0.2\headwidth-2\fboxsep][c]{\strut\textcolor{ehAzul1}{\bfseries\thepage}}}}
	\renewcommand{\headrulewidth}{0pt}
}
\makeatother

% +-----------------+


% +-----------------+
% |  Bibliografías  |
% +-----------------+

% Librerías
\usepackage[sectionbib]{natbib}			% Librería para tener una bibliografía por capítulo
\usepackage{chapterbib}					% Hmmm, esta librería parece hacer lo que creo que hace la anterior, pero ambas son necesarias
\usepackage{apalike}					% Librería para usar APA en los artículos

% Comando para poner bien las referencias en cada artículo
\newcommand{\referencias}{
	\renewcommand{\bibname}{\normalsize Referencias\vspace{-1.5cm}}		% Cambia el nombre a Referencias y reduce el espacio entre este y la bibliografía
	\bibliographystyle{apalike}											% Usar el estilo APA
	{\footnotesize \bibliography{tex/bibliografia-No1}}					% Especifica el archibo que contiene la información bibliográfica
}

% +-----------------+


% +----------------------+
% |  Info del documento  |
% +----------------------+

\title{\textit{\textcolor{ehAzul1}{\textbf{ex}}\textcolor{ehAzul2}{$\rightarrow$}\textcolor{ehAzul1}{\vspace{-0.4cm}hypothesi}}}		 				% Título de la revista
\date{}								% No poner ninguna fecha para que no salga

% +-----------------+


% +---------------+
% |  "Variables"  |
% +---------------+

% Los que más cambian
\newcommand{\numeroRevista}{1}										% El numero de la revista
\newcommand{\semestre}{2019-1}										% El semestre en el que se hizo el número
\newcommand{\representante}{Martín Buenahora Bonilla}				% Representante de estudiantes actual
\newcommand{\editores}{	Pablo Rivas Robledo\tabularnewline			% Miembros del comité editorial
						& Martín Buenahora Bonilla}
\newcommand{\profesores}{	Juan Camilo Espejo Serna\tabularnewline
							& Mariano Lozano Ramírez (corrección de estilo)}

% Estos, en principio, no deberían cambiar tanto
\newcommand{\rector}{Obdulio Velásquez Posada}				% Rector/a actual de la U
\newcommand{\decano}{Bogdan Piotrowski}						% Decano/a actual de la facultad
\newcommand{\director}{Carmen Elena Arboleda}				% Director/a actual del programa

% +-----------------+


\begin{document}

% +-----------------------+
% |  Preambula articulii  |
% +-----------------------+

\maketitle									% Título del libro/revista

{\small

\begin{center}
	{\Huge\textit{\textcolor{ehAzul1}{\textbf{ex}}\textcolor{ehAzul2}{$\rightarrow$}\textcolor{ehAzul1}{\vspace{-0.4cm}hypothesi}}}\\[0.5cm]
	Revista de estudiantes de Filosofía -- Una publicación de
	\includegraphics[width=9.021cm,height=1.984cm]{img/logo-neuronum.png}
\end{center}

\begin{longtable}[]{@{}ll@{}}
	\toprule[0.5pt]
	\midrule[0.5pt]
	\endhead
%	Rector & \rector{}\tabularnewline
%	&\tabularnewline
%	Decano Facultad de Filosofía y Ciencias Humanas & \decano{}\tabularnewline
%	&\tabularnewline
%	Directora del Programa de Filosofía & \director{}\tabularnewline
%	&\tabularnewline
	Representante de Estudiantes & \representante{}\tabularnewline
	&\tabularnewline
	Profesor acompañante & \profesores{}\tabularnewline
	&\tabularnewline
	Edición a cargo de & \editores{}\tabularnewline
	\bottomrule[0.5pt]
\end{longtable}
%\vspace{-2.25em}\hspace{18.7em}\textcolor{white}{La mónada de Leibniz (Dirección espiritual)}
\newpage

\begin{center}
	\includegraphics[width=2.24419in,height=0.78333in]{img/CC-BY-NC-ND.png}
\end{center}

\begin{flushleft}

El material elaborado para esta publicación puede ser\\
distribuido, copiado y exhibido por terceros para\\
fines académicos, siempre y cuando se cite la fuente.\\
No se puede obtener ningún beneficio comercial y las\\
obras derivadas se deben presentar bajo los mismos\\
términos de licencia que el trabajo original.

El contenido de los artículos es de exclusiva\\
responsabilidad de los autores. La Universidad de La\\
Sabana no se compromete con las opiniones que se\\
expresen.

Semestre \semestre{}. Número \numeroRevista{}, Volumen \volumen{}.

\textit{ex hypothesi} es una revista de publicación anual,\\
editada por estudiantes de Filosofía de la Universidad\\
de La Sabana.

Revista \textit{ex hypothesi} – Revista de estudiantes de\\
filosofía – Una publicación de Neuronum.\\
Neurociencias al servicio de la educación, para el\\
desarrollo personal y social.

Campus del Puente del Común

Km 7 Autopista Norte de Bogotá

Chía, Cundinamarca, Colombia

\noindent Contacto: \href{mailto:exhypothesi@unisabana.edu.co}{exhypothesi@unisabana.edu.co}

\end{flushleft}

\newpage

%\begin{flushleft}
%	\textit{ex hypothesi} agradece a todas las\\
%	personas que contribuyeron de una u\\
%	otra manera con la elaboración de este\\
%	número. A ellos les agradecemos\\
%	inmensamente por sus aportes.\\[6cm]
%\end{flushleft}

%\begin{flushright}
%	\textit{ex hypothesi}\\
%	/ˌ\textepsilon ks h\textturnv \textsci ˈp\textturnscripta \texttheta \textschwa s\textturnv \textsci/\\
%	De acuerdo con la hipótesis propuesta
%\end{flushright}

{\parindent0pt
(...) que ninguna de nuestras acciones\\
se olvida; que todo entra en la cuenta,\\
hasta las palabras más ociosas y hasta\\
una cucharada de agua bien empleada;\\
en fin, que todo tiene que resultar para\\
el mayor bien de los buenos, que los\\
justos serán como soles, y que ni\\
nuestros sentidos ni nuestro espíritu\\
han gustado nunca nada que se\\
aproxime a la felicidad que Dios\\
prepara a los que lo aman.

Leibniz, Discurso de Metafísica, § 37
}			% Fin del \parindent0pt

}			% Fin del \small				% Información sobre la revista y los derechos de autor

\tableofcontents							% Índice

\small					% La tipografía normal del documento es
						% small, no normalzise. Esto se debe a que
						% era la forma más sencilla de conseguir 
						% _este_ tamaño de letra. 10pt, 11pt, 12pt
						% etc. no daban este resultado, ya que siempre
						% era o muy grande o muy pequeño. Aunque esto
						% seguramente va a cambiar pronto, pero me
						% parece bueno recordar por qué carajos
						% hay un  small aquí.

% Configurar pié de página con la info de la revista. Esto está aquí para que aparezca _después_ y no antes del índice, que se vería feo
\fancyfoot[LE,LO]{{\footnotesize\textit{ex hypothesi}\\Una publicación de \textit{Neuronum}\\\vspace{-0.2cm}ISSN: 2422-5193}}
\renewcommand{\footrulewidth}{0.5pt}
\futurelet\TMPfootrule\def\footrule{{\color{ehAzul1}\TMPfootrule}}

\onehalfspacing								% Interlineado 1.5

% +-----------------+


% +-------------+
% |  Artículos  |
% +-------------+

%\include{./tex/}

% +-----------------+


\end{document}
